\par\indent\SourceLeaf{225}%
III. Angulus inter duo plana aequalis est angulo sphaerico inter circulos maximos illa repraesentantes, et proin etiam per arcum inter horum circulorum maximorum polos interceptum mensuratur. Et perinde inclinatio rectae ad planum mensuratur per arcum, a puncto, quod respondet directioni rectae, ad circulum maximum, qui plani situm repraesentat, normaliter ductum.

IV. Denotantibus \(x,y,z;\ x',y',z'\) coordinatas duorum punctorum, \(r\) eorundem distantiam, atque \(L\) punctum, quod in superficie sphaerica repraesentat directionem rectae a puncto priore ad posterius ductae, erit
\[
x'=x+r\cos(1)L,\qquad y'=y+r\cos(2)L,\qquad z'=z+r\cos(3)L
\]
\par\indent V. Hinc facile sequitur, haberi generaliter
\[
\cos(1)L^2+\cos(2)L^2+\cos(3)L^2=1
\]
\noindent nec non, denotante \(L'\) quodcunque aliud punctum superficiei sphaericae, esse
\[
\cos(1)L\mathord{.}\cos(1)L'+\cos(2)L\mathord{.}\cos(2)L'+\cos(3)L\mathord{.}\cos(3)L'=\cos LL'
\]
\par\indent VI. \textsc{Theorema.} \textit{Denotantibus \(L,L',L'',L'''\) quatuor puncta in superficie sphaerae, atque \(A\) angulum, quem arcus \(LL',L''L'''\) in puncto concursus sui formant, erit}
\[
\cos LL''\mathord{.}\cos L'L'''-\cos LL'''\mathord{.}\cos L'L''=\sin LL'\mathord{.}\sin L''L'''\mathord{.}\cos A
\]
\par\indent\textit{Dem.} Denotet litera \(A\) insuper punctum concursus ipsum, statuaturque
\[
AL=t,\quad AL'=t',\quad AL''=t'',\quad AL'''=t'''
\]
\noindent Habemus itaque:
\[
\begin{aligned}
\cos LL''&=\cos t\cos t''+\sin t\sin t''\cos A\\
\cos L'L'''&=\cos t'\cos t'''+\sin t'\sin t'''\cos A\\
\cos LL'''&=\cos t\cos t'''+\sin t\sin t'''\cos A\\
\cos L'L''&=\cos t'\cos t''+\sin t'\sin t''\cos A
\end{aligned}
\]
\noindent et proin
\[
\begin{aligned}
&\cos LL''\mathord{.}\cos L'L'''-\cos LL'''\mathord{.}\cos L'L''\\
&\quad=\cos A(\cos t\cos t''\sin t'\sin t'''+\cos t'\cos t'''\sin t\sin t''\\
&\qquad\qquad-\cos t\cos t'''\sin t'\sin t''-\cos t'\cos t''\sin t\sin t''')\\
&\quad=\cos A(\cos t\sin t'-\sin t\cos t')(\cos t''\sin t'''-\sin t''\cos t''')\\
&\quad=\cos A\mathord{.}\sin(t'-t)\mathord{.}\sin(t'''-t'')\\
&\quad=\cos A\mathord{.}\sin LL'\mathord{.}\sin L''L'''
\end{aligned}
\]
